\documentclass[aspectratio=169]{beamer}

\usepackage{amsmath,amssymb}
\usepackage{graphicx}
\usepackage{tikz}
\usepackage{listings}
\usetikzlibrary{positioning,calc,arrows.meta,fit}

% =========================
% Colors
% =========================
\definecolor{PKUDarkPurple}{RGB}{63,63,191}
\definecolor{PKULightPurple}{RGB}{170,170,225}
\definecolor{PKUGrayBg}{RGB}{255,255,255}
\definecolor{PKUTextBlue}{RGB}{45,45,190}
\definecolor{CodeGray}{RGB}{245,245,245}
\definecolor{dkgreen}{rgb}{0,0.5,0}
\definecolor{gray}{rgb}{0.4,0.4,0.4}
\definecolor{mauve}{rgb}{0.58,0,0.82}

% =========================
% Global style
% =========================
\setbeamercolor{background canvas}{bg=PKUGrayBg}
\setbeamercolor{normal text}{fg=black,bg=PKUGrayBg}
\setbeamercolor{frametitle}{fg=PKUTextBlue,bg=PKUGrayBg}

\setbeamerfont{normal text}{size=\fontsize{15pt}{18pt}\selectfont}
\setbeamerfont{frametitle}{size=\fontsize{15pt}{19pt}\selectfont,series=\bfseries}
\setbeamerfont{itemize/enumerate body}{size=\fontsize{12pt}{15pt}\selectfont}
\setbeamerfont{itemize/enumerate subbody}{size=\fontsize{11pt}{14pt}\selectfont}
\setbeamerfont{footline}{size=\fontsize{10pt}{12pt}\selectfont}

\setbeamercolor{itemize item}{fg=PKUDarkPurple}
\setbeamertemplate{navigation symbols}{}
\setbeamertemplate{items}[ball]

\setbeamertemplate{headline}{}
\setbeamertemplate{footline}{}

% =========================
% Code style
% =========================
\lstset{
  basicstyle=\ttfamily\footnotesize,
  backgroundcolor=\color{CodeGray},
  keywordstyle=\color{blue},
  commentstyle=\color{dkgreen},
  stringstyle=\color{mauve},
  showstringspaces=false,
  breaklines=true,
  frame=single,
  columns=fullflexible
}

% =========================
% Custom headline
% =========================
\newcommand{\makecustomheadline}{
\begin{tikzpicture}[remember picture,overlay]
    \fill[PKUDarkPurple]
        (current page.north west)
        rectangle
        ([xshift=.50\paperwidth,yshift=-0.35cm]current page.north west);

    \fill[PKULightPurple]
        ([xshift=.50\paperwidth]current page.north west)
        rectangle
        ([xshift=\paperwidth,yshift=-0.35cm]current page.north west);
\end{tikzpicture}
}

% =========================
% Custom footline
% =========================
\newcommand{\makecustomfootline}{
\begin{tikzpicture}[remember picture,overlay]
    \fill[PKUDarkPurple]
        (current page.south west)
        rectangle
        ([xshift=.33\paperwidth,yshift=0.42cm]current page.south west);

    \fill[PKULightPurple]
        ([xshift=.33\paperwidth]current page.south west)
        rectangle
        ([xshift=.67\paperwidth,yshift=0.42cm]current page.south west);

    \fill[PKULightPurple]
        ([xshift=.67\paperwidth]current page.south west)
        rectangle
        ([xshift=\paperwidth,yshift=0.42cm]current page.south west);

    \node[anchor=center, text=white, font=\small]
        at ([xshift=.165\paperwidth,yshift=0.21cm]current page.south west)
        {\textrm{huangxun@pku.edu.cn}};

    \node[anchor=center, text=white, font=\small]
        at ([xshift=.50\paperwidth,yshift=0.21cm]current page.south west)
        {AI Enabled Control Engineering};

    \node[anchor=center, text=black, font=\small]
        at ([xshift=.83\paperwidth,yshift=0.21cm]current page.south west)
        {2026 Summer};

    \node[anchor=center, text=black, font=\small]
        at ([xshift=.965\paperwidth,yshift=0.21cm]current page.south west)
        {\insertframenumber/\inserttotalframenumber};
\end{tikzpicture}
}

\addtobeamertemplate{background}{\makecustomheadline\makecustomfootline}{}

% =========================
% Helpers
% =========================
\newcommand{\smallblue}[1]{\textcolor{PKUTextBlue}{#1}}
\newcommand{\mat}[1]{\mathbf{#1}}

% =========================
% Title info
% =========================
\title{\textcolor{PKUTextBlue}{AI Enabled Control Engineering}}
\author{\textcolor{blue}{Professor Xun Huang}}
\institute{
Aeronautics and Astronautics\\
College of Engineering\\
Peking University\\[0.25cm]
huangxun@pku.edu.cn\\
https://xunger99.github.io/xunger/
}
\date{}

\setbeamertemplate{title page}{
    \vspace{1.0cm}
    \begin{center}
        {\fontsize{18}{22}\selectfont\bfseries \inserttitle\par}
        \vspace{0.3cm}
        {\fontsize{15}{18}\selectfont Lecture 8: Optimal Control Experiments on the Rotary Inverted Pendulum\par}
        \vspace{0.3cm}
        {\fontsize{15}{18}\selectfont \insertauthor\par}
        \vspace{0.15cm}
        {\fontsize{11}{13}\selectfont TA: Zhixiang Ju \quad Haozhe Wang\par}
        \vspace{0.4cm}
        {\fontsize{8}{11}\selectfont \insertinstitute\par}
    \end{center}
}

\begin{document}

% ---------------------------------
\begin{frame}
\titlepage
\end{frame}

% ---------------------------------
\begin{frame}{Recap: from Class 7 to Class 8}
\begin{itemize}
    \item In Class 7, we assembled the real rotary inverted pendulum circuit.
    \vspace{0.25cm}
    \item We tested the potentiometer, encoder, DC motor, and data logging pipeline.
    \vspace{0.25cm}
    \item We connected Bellman's optimality principle to HJB, LQR, MPC, and Q-learning.
    \vspace{0.25cm}
    \item Today we use model-based optimal control to stabilize the real hardware.
\end{itemize}
\end{frame}

% ---------------------------------
\begin{frame}{Lecture 8}
\begin{itemize}
    \item Review the local linear model of the rotary inverted pendulum.
    \vspace{0.28cm}
    \item Derive the discrete infinite-horizon LQR controller.
    \vspace{0.28cm}
    \item Compute the LQR gain and deploy it to the real STM32 system.
    \vspace{0.28cm}
    \item Introduce linear MPC as constrained finite-horizon optimal control.
    \vspace{0.28cm}
    \item Run the MPC real-hardware experiment and compare it with LQR.
\end{itemize}
\end{frame}

% ---------------------------------
\begin{frame}{Learning objectives}
By the end of this lab, students should be able to:
\vspace{0.25cm}
\begin{itemize}
    \item explain why LQR and MPC are local stabilizing controllers for the upright region;
    \vspace{0.2cm}
    \item compute a discrete-time LQR feedback gain from matrices $A_d,B_d,Q,R$;
    \vspace{0.2cm}
    \item understand how low-pass filtered angle differences provide angular-velocity estimates;
    \vspace{0.2cm}
    \item modify the firmware constants and run a real LQR balancing experiment;
    \vspace{0.2cm}
    \item understand how MPC adds finite-horizon prediction and PWM constraints.
\end{itemize}
\end{frame}

% ---------------------------------
\begin{frame}{Experimental simplification}
\begin{alertblock}{No swing-up in this class}
\small
In this experiment, we do not use DQN and we do not use an automatic swing-up stage.
\end{alertblock}

\vspace{0.25cm}

\begin{itemize}
    \item The controller is only valid near the upright equilibrium.
    \vspace{0.25cm}
    \item Students should gently hold the pendulum close to the upright position.
    \vspace{0.25cm}
    \item Then run the Python script and press Enter in the terminal to send the \texttt{GO} command.
    \vspace{0.25cm}
    \item After \texttt{GO}, STM32 records the startup arm position as $\theta=0$ and starts closed-loop control.
\end{itemize}
\end{frame}

% ---------------------------------
\begin{frame}{From continuous model to sampled-data model}

The physical rotary inverted pendulum is a continuous-time system:
\[
\dot{x}(t)=Ax(t)+Bu(t).
\]



However, the STM32 controller does not update the motor command continuously.



In our experiment:
\[
T_s=0.005\ \mathrm{s},
\qquad
f_s=200\ \mathrm{Hz}.
\]



During one sampling interval, the PWM command is held constant:
\[
u(t)=u_k,
\qquad
t\in[kT_s,(k+1)T_s).
\]



\begin{block}{Key point}
\small
The controller runs in discrete time. Therefore, the LQR and MPC gains used on STM32 should be computed from the discrete model, not directly from the continuous model.
\end{block}

\end{frame}




% ---------------------------------
\begin{frame}{Zero-order-hold discretization}

For the continuous-time linear model
\[
\dot{x}(t)=Ax(t)+Bu(t),
\]
with zero-order hold input \(u(t)=u_k\), the exact sampled model is
\[
x_{k+1}=A_d x_k+B_d u_k.
\]



The matrices are
\[
A_d=e^{AT_s},
\]
\[
B_d=\int_0^{T_s}e^{A\tau}B\,d\tau.
\]



In Python, students can compute them by
\[
(A_d,B_d)=\texttt{cont2discrete}(A,B,T_s).
\]





\end{frame}





% ---------------------------------

% ---------------------------------
\begin{frame}{How do we get angular velocities?}
The sensors directly measure only two angles:
\[
\theta_k,
\qquad
\alpha_k.
\]

\vspace{0.2cm}

In this class, use angle difference with low-pass filtering:
\[
\dot\theta^{\mathrm{raw}}_k
=
\frac{\theta_k-\theta_{k-1}}{T_s},
\qquad
\dot\alpha^{\mathrm{raw}}_k
=
\frac{\mathrm{wrap}(\alpha_k-\alpha_{k-1})}{T_s}.
\]

\vspace{0.2cm}

Then apply first-order low-pass filtering:
\[
\dot\theta_k=(1-\lambda)\dot\theta_{k-1}+\lambda\dot\theta^{\mathrm{raw}}_k,
\]
\[
\dot\alpha_k=(1-\lambda)\dot\alpha_{k-1}+\lambda\dot\alpha^{\mathrm{raw}}_k.
\]
\end{frame}

% ---------------------------------
\begin{frame}{Why low-pass filtering is necessary}
\begin{itemize}
    \item Direct difference amplifies sensor noise.
    \vspace{0.25cm}
    \item Potentiometer quantization and encoder count jumps may produce noisy velocity estimates.
    \vspace{0.25cm}
    \item LQR and MPC both use the full state $x=[\theta,\dot\theta,\alpha,\dot\alpha]^\top$.
    \vspace{0.25cm}
    \item Therefore, velocity estimates must be smooth enough for real-time feedback.
\end{itemize}

\vspace{0.25cm}
\begin{block}{Practical choice}
\small
The firmware uses a fixed low-pass coefficient. Students can compare different filtering strengths after the basic experiment works.
\end{block}
\end{frame}

% ---------------------------------
\begin{frame}{LQR problem setup}
For the discrete linear system
\[
x_{k+1}=A_d x_k+B_d u_k,
\]
choose the infinite-horizon quadratic cost
\[
J=\sum_{k=0}^{\infty}
\left(
x_k^\top Qx_k+u_k^\top Ru_k
\right).
\]

\vspace{0.2cm}

Here
\[
Q\succeq 0,
\qquad
R\succ0.
\]

\vspace{0.2cm}
\begin{block}{Physical meaning}
\small
$Q$ penalizes state deviation from the upright equilibrium. $R$ penalizes large PWM commands.
\end{block}
\end{frame}

% ---------------------------------
\begin{frame}{Choosing $Q$ and $R$ for the RIP}
The state order is
\[
x=\begin{bmatrix}\theta&\dot\theta&\alpha&\dot\alpha\end{bmatrix}^\top.
\]

A typical choice is
\[
Q=\mathrm{diag}(1,\ 0.05,
80,\ 2),
\qquad
R=0.001.
\]

\vspace{0.2cm}
\begin{itemize}
    \item $q_\alpha$ is large because the pendulum angle is the most critical state.
    \vspace{0.18cm}
    \item $q_{\dot\alpha}$ suppresses pendulum angular velocity.
    \vspace{0.18cm}
    \item $q_\theta$ limits long-term arm drift.
    \vspace{0.18cm}
    \item $R$ controls how aggressively the motor is used.
\end{itemize}
\end{frame}

% ---------------------------------
\begin{frame}{Discrete LQR value function}
Assume the optimal value function has a quadratic form:
\[
V(x_k)=x_k^\top P x_k,
\qquad
P=P^\top\succeq0.
\]

Bellman's optimality principle gives
\[
V(x_k)=\min_{u_k}
\left\{
x_k^\top Qx_k+u_k^\top Ru_k+V(x_{k+1})
\right\}.
\]

\vspace{0.15cm}

Using
\[
x_{k+1}=A_dx_k+B_du_k,
\]
we obtain
\[
x_k^\top P x_k
=\min_{u_k}
\left\{
x_k^\top Qx_k+u_k^\top Ru_k
+(A_dx_k+B_du_k)^\top P(A_dx_k+B_du_k)
\right\}.
\]
\end{frame}

% Fix for beamer parsing: define empty macro used above
% ---------------------------------
\begin{frame}{LQR feedback gain}
The terms depending on $u_k$ are quadratic:
\[
u_k^\top(R+B_d^\top PB_d)u_k
+2u_k^\top B_d^\top PA_dx_k.
\]

\vspace{0.2cm}

Taking derivative with respect to $u_k$ and setting it to zero gives
\[
(R+B_d^\top PB_d)u_k+B_d^\top PA_dx_k=0.
\]

\vspace{0.2cm}

Thus
\[
u_k^*=-Kx_k,
\]
where
\[
K=(R+B_d^\top PB_d)^{-1}B_d^\top PA_d.
\]
\end{frame}

% ---------------------------------
\begin{frame}{Discrete algebraic Riccati equation}
The matrix $P$ satisfies the discrete algebraic Riccati equation:
\[
P
=
A_d^\top P A_d
-
A_d^\top P B_d
(R+B_d^\top P B_d)^{-1}
B_d^\top P A_d
+Q.
\]

\vspace{0.2cm}

After solving for $P$, compute
\[
K=(R+B_d^\top PB_d)^{-1}B_d^\top PA_d.
\]

\vspace{0.2cm}
\begin{alertblock}{For this lab}
\small
We use infinite-horizon discrete LQR. Therefore, we solve the algebraic Riccati equation and do not need a terminal condition $P(T)=H$.
\end{alertblock}
\end{frame}

% ---------------------------------
\begin{frame}[fragile]{Student task: compute the LQR gain}
Students should find or load the model matrices $A_d$ and $B_d$ from the provided code or handout.

\vspace{0.15cm}

Then compute the gain in Python:
\begin{lstlisting}[basicstyle=\ttfamily\tiny]
import numpy as np
from scipy.linalg import solve_discrete_are

Q = np.diag([1.0, 0.05, 80.0, 2.0])
R = np.array([[0.001]])

P = solve_discrete_are(Ad, Bd, Q, R)
K = np.linalg.inv(R + Bd.T @ P @ Bd) @ (Bd.T @ P @ Ad)
print(K)
\end{lstlisting}

\vspace{0.12cm}
\begin{block}{Output}
\small
The result is a row vector $K=[k_\theta,k_{\dot\theta},k_\alpha,k_{\dot\alpha}]$.
\end{block}
\end{frame}

% ---------------------------------
\begin{frame}{How to put $K$ into firmware}
The control law is
\[
u_k=-Kx_k.
\]

For our firmware, insert the four gains into:
\[
K_\theta,
\quad
K_{\dot\theta},
\quad
K_\alpha,
\quad
K_{\dot\alpha}.
\]

\vspace{0.2cm}

Then the firmware computes
\[
u_{
m raw}
=-(K_\theta\theta
+K_{\dot\theta}\dot\theta
+K_\alpha\alpha
+K_{\dot\alpha}\dot\alpha).
\]

\vspace{0.2cm}

Finally, the PWM command is saturated:
\[
u=\mathrm{sat}(u_{
m raw},-u_{\max},u_{\max}).
\]
\end{frame}

% ---------------------------------
\begin{frame}{LQR hardware experiment: preparation}

\begin{enumerate}
    \item Upload the LQR firmware to STM32.
    \vspace{0.18cm}

    \item Calibrate the potentiometer zero position:
    hold the pendulum upright, read the current raw value, and update \texttt{POT\_ZERO} in the firmware.
    \vspace{0.18cm}

    \item Upload the firmware again after modifying \texttt{POT\_ZERO}.
    \vspace{0.18cm}

    \item Close the Arduino Serial Monitor.
    \vspace{0.18cm}

    \item Connect motor power only after the wiring is checked.
    \vspace{0.18cm}

    \item Hold the pendulum gently near the upright position before pressing Enter in the Python terminal.
    \vspace{0.18cm}

    \item Keep fingers clear of the rotating arm.
\end{enumerate}

\vspace{0.15cm}


\end{frame}

% ---------------------------------
\begin{frame}{LQR hardware experiment: run}
\begin{enumerate}
    \item Open a terminal in the LQR experiment folder.
    \vspace{0.25cm}
    \item Run the Python script:
    \vspace{0.1cm}
    \begin{center}
    \texttt{python rip\_lqr\_run.py}
    \end{center}
    \vspace{0.2cm}
    \item Hold the pendulum close to upright.
    \vspace{0.25cm}
    \item Press Enter in the Python terminal to send \texttt{GO}.
    \vspace{0.25cm}
    \item The script records data and saves a CSV log and figures.
\end{enumerate}
\end{frame}

% ---------------------------------
% ---------------------------------
\begin{frame}{LQR hardware experiment: expected data}

The Python script records:
\[
t,\quad \theta,\quad \alpha,
\quad \dot\theta_{\mathrm{LPF}},
\quad \dot\alpha_{\mathrm{LPF}},
\quad u.
\]

\vspace{0.2cm}

Here, the angular velocities are not directly measured.
They are computed from angle differences and then filtered:
\[
\dot\theta_{\mathrm{raw}}(k)
=
\frac{\theta_k-\theta_{k-1}}{T_s},
\qquad
\dot\alpha_{\mathrm{raw}}(k)
=
\frac{\alpha_k-\alpha_{k-1}}{T_s}.
\]

\[
\dot{x}_{\mathrm{LPF}}(k)
=
(1-\lambda)\dot{x}_{\mathrm{LPF}}(k-1)
+
\lambda \dot{x}_{\mathrm{raw}}(k).
\]

\vspace{0.15cm}

The generated plots should include:
\begin{itemize}
    \item angle response: $\theta(t)$ and $\alpha(t)$;
    \item filtered angular velocity response: $\dot\theta_{\mathrm{LPF}}(t)$ and $\dot\alpha_{\mathrm{LPF}}(t)$;
    \item PWM command $u(t)$.
\end{itemize}

\vspace{0.15cm}


\end{frame}

% ---------------------------------
\begin{frame}{Why introduce MPC after LQR?}
LQR gives an explicit feedback law:
\[
u=-Kx.
\]

\vspace{0.15cm}

This is simple and fast, but standard LQR does not explicitly handle input constraints.

\vspace{0.15cm}

In real hardware, the PWM command must satisfy
\[
-u_{\max}\le u\le u_{\max}.
\]

\vspace{0.2cm}
\begin{block}{MPC idea}
\small
MPC solves a finite-horizon optimal control problem at each sampling instant and includes PWM constraints directly inside the optimization problem.
\end{block}
\end{frame}

% ---------------------------------
\begin{frame}{MPC prediction model}
Use the same discrete linear model:
\[
x_{k+1}=A_dx_k+B_du_k.
\]

At current time $k$, predict future states:
\[
x_{k+i+1|k}
=A_dx_{k+i|k}+B_du_{k+i|k},
\qquad i=0,1,\ldots,N-1.
\]

The initial predicted state is
\[
x_{k|k}=x_k.
\]

\vspace{0.15cm}
\begin{block}{Prediction horizon}
\small
In the provided embedded MPC code, the prediction horizon is small enough for STM32 real-time implementation at 200 Hz.
\end{block}
\end{frame}

% ---------------------------------
\begin{frame}{Stacked prediction form}
Define the future input sequence
\[
U=\begin{bmatrix}
u_{k|k}&u_{k+1|k}&\cdots&u_{k+N-1|k}
\end{bmatrix}^\top.
\]

Stack future predicted states:
\[
X=\begin{bmatrix}
x_{k+1|k}^\top&x_{k+2|k}^\top&\cdots&x_{k+N|k}^\top
\end{bmatrix}^\top.
\]

Then
\[
X=S_xx_k+S_uU.
\]

\vspace{0.15cm}
\begin{block}{Meaning}
\small
Once the current state $x_k$ is known, all future predicted states are linear functions of the future input sequence $U$.
\end{block}
\end{frame}

% ---------------------------------
\begin{frame}{MPC cost function}
The finite-horizon MPC cost is
\[
J
=
\sum_{i=0}^{N-1}
\left(
x_{k+i|k}^\top Qx_{k+i|k}
+
u_{k+i|k}^\top Ru_{k+i|k}
\right)
+x_{k+N|k}^\top P x_{k+N|k}.
\]

\vspace{0.15cm}

Use the same stage weights as LQR:
\[
Q=\mathrm{diag}(1,0.05,80,2),
\qquad
R=0.001.
\]

\vspace{0.15cm}
\begin{block}{Terminal weight}
\small
The terminal weight $P$ can be chosen as the infinite-horizon LQR Riccati solution $P_\infty$ to approximate the tail cost beyond the prediction window.
\end{block}
\end{frame}

% ---------------------------------
\begin{frame}{From MPC cost to quadratic programming}
Using
\[
X=S_xx_k+S_uU,
\]
the part of the cost depending on $U$ can be written as
\[
J_U(U)=\frac12U^\top HU+f^\top U,
\]
where
\[
f=Gx_k.
\]

\vspace{0.2cm}

The matrices $H$ and $G$ depend only on
\[
A_d,
B_d,
Q,
R,
P,
N.
\]

\vspace{0.15cm}
\begin{block}{Embedded implementation}
\small
$H$ and $G$ are computed offline and stored in the firmware. Online computation only updates $f=Gx_k$.
\end{block}
\end{frame}

% ---------------------------------
\begin{frame}{MPC input constraint}
The motor PWM command is bounded:
\[
-u_{\max}\le u_{k+i|k}\le u_{\max},
\qquad i=0,1,\ldots,N-1.
\]

Therefore, the MPC problem is
\[
\min_U\ \frac12U^\top HU+(Gx_k)^\top U,
\]
\[
\mathrm{s.t.}\quad
-u_{\max}\le U_i\le u_{\max}.
\]

\vspace{0.15cm}
\begin{block}{Main difference from LQR}
\small
LQR computes an unconstrained feedback and then the firmware may saturate it. MPC considers the PWM bound inside the optimization problem.
\end{block}
\end{frame}

% ---------------------------------
\begin{frame}{Projected-gradient solver on STM32}
The gradient of the QP objective is
\[
\nabla J(U)=HU+f.
\]

A projected-gradient iteration is
\[
U^{(j+1)}=
\Pi_{[-u_{\max},u_{\max}]}
\left(
U^{(j)}-\eta(HU^{(j)}+f)
\right).
\]

\vspace{0.15cm}

Here $\Pi$ means element-wise saturation:
\[
U_i\leftarrow \max(-u_{\max},\min(u_{\max},U_i)).
\]

\vspace{0.15cm}
\begin{block}{Real-time strategy}
\small
Use a fixed number of iterations and warm start from the previous control sequence.
\end{block}
\end{frame}

% ---------------------------------
\begin{frame}{MPC rolling optimization}
At each control cycle:
\begin{enumerate}
    \item read the current state $x_k$;
    \vspace{0.15cm}
    \item compute $f=Gx_k$;
    \vspace{0.15cm}
    \item warm start the control sequence $U$;
    \vspace{0.15cm}
    \item run projected-gradient iterations;
    \vspace{0.15cm}
    \item apply only the first input:
    \[
    u_k=u^*_{k|k};
    \]
    \vspace{0.05cm}
    \item repeat at the next 5 ms control cycle.
\end{enumerate}
\end{frame}

% ---------------------------------
\begin{frame}{MPC hardware experiment}
\begin{enumerate}
    \item Upload the MPC firmware to STM32.
    \vspace{0.22cm}
    \item Close the Arduino Serial Monitor.
    \vspace{0.22cm}
    \item Hold the pendulum near the upright position.
    \vspace{0.22cm}
    \item Run the Python script:
    \vspace{0.08cm}
    \begin{center}
    \texttt{python rip\_mpc\_run.py}
    \end{center}
    \vspace{0.15cm}
    \item Press Enter in the terminal to send \texttt{GO}.
    \vspace{0.22cm}
    \item Record angle, angular velocity, PWM, and computation-time logs.
\end{enumerate}
\end{frame}

% ---------------------------------
\begin{frame}{Comparing LQR and MPC in the lab}
\begin{center}
\begin{tabular}{c|c|c}
\hline
Item & LQR & MPC \\
\hline
Control law & explicit $u=-Kx$ & QP-defined implicit feedback \\
Constraint handling & saturation after control & constraint inside optimization \\
Online computation & very small & projected-gradient iterations \\
Tuning parameters & $Q,R$ & $Q,R,P,N,u_{\max}$ \\
Hardware focus & fast local stabilization & constrained local stabilization \\
\hline
\end{tabular}
\end{center}

\vspace{0.25cm}
\begin{block}{Student comparison task}
\small
Compare angle curves, PWM distributions, and whether the controller frequently hits the PWM bound.
\end{block}
\end{frame}

% ---------------------------------
% ---------------------------------
\begin{frame}{Lab report requirements}

Each group should record:
\begin{itemize}
    \item the model matrices or file location used for computing \(K\);


    \item at least three different LQR weight settings:
    \[
    Q=\mathrm{diag}(q_\theta,q_{\dot\theta},q_\alpha,q_{\dot\alpha}),
    \qquad R=r;
    \]
   

    \item the computed LQR gain \(K\) for each \(Q,R\) setting;
   

    \item one successful experimental log for each LQR setting;
   

    \item one successful MPC experimental log;
  

    \item comparison plots of \(\theta,\alpha,\dot\theta_{\mathrm{LPF}},\dot\alpha_{\mathrm{LPF}}\) and PWM for different LQR settings;
  

    \item a short comparison of LQR tuning effects and MPC performance.
\end{itemize}

\end{frame}



% ---------------------------------
\begin{frame}{Summary of Lecture 8}
\begin{itemize}
    \item LQR and MPC both rely on the local discrete model $x_{k+1}=A_dx_k+B_du_k$.

    \item LQR solves a discrete Riccati equation and gives a fast explicit feedback law.

    \item MPC solves a finite-horizon constrained optimization problem at each sampling time.

    \item In the real experiment, students manually place the pendulum near upright and press Enter in Python to start control.

    \item Low-pass filtered angle differences provide the velocity estimates needed by both controllers.

    \item The key experimental comparison is between unconstrained feedback plus saturation and constrained rolling optimization.
\end{itemize}
\end{frame}

\end{document}
